\documentclass[11pt]{article}
\usepackage[T1]{fontenc}
\usepackage{amsmath}
\usepackage{booktabs}

\title{{{title}}}
\author{{{author}}}
\date{}

\begin{document}
\maketitle

\begin{abstract}
This is a placeholder abstract. Replace this text with a summary of your research, findings, or main argument.
\end{abstract}

\section{Introduction}

This is the introduction section. Present the background and context for your work. You can cite other authors, such as the foundational work by Smith \cite{smith2020}.

\section{Methods}

Describe your methodology or approach in this section. Include relevant details about your process, tools, or theoretical framework.

\section{Results}

Present your findings in this section. The equation below shows an important relationship:
\begin{equation}
  y = mx + b
  \label{eq:linear}
\end{equation}

The results can be summarized in a table, as shown in Table~\ref{tbl:example}.

\begin{table}[ht]
  \centering
  \caption{Example results table.}
  \label{tbl:example}
  \begin{tabular}{lcc}
    \toprule
    Variable & Value & Unit \\
    \midrule
    Sample size & 100 & \\
    Mean & 42.5 & units \\
    Std Dev & 5.2 & units \\
    \bottomrule
  \end{tabular}
\end{table}

\section{Discussion}

Interpret your results in the context of existing literature. Discuss the implications, limitations, and potential future directions for this work.

\bibliographystyle{plain}
\bibliography{references}

\end{document}
